\chapter{System Requirements and Tiered Architecture}
\label{ch:architecture}

This chapter converts the research questions into system requirements and
defines the boundary between the stationary perception workflow and the
two-vehicle mission software. It presents intended behaviour and verified
implementation status; measurements of performance are reserved for
Chapters~6 and 7.

\section{Functional and non-functional requirements}
\label{sec:architecture:requirements}

The requirements were derived from the research questions, the available
hardware and the need to preserve auditable evidence. A requirement can be
implemented without being experimentally validated. Table~\ref{tab:requirements}
therefore identifies both its architectural purpose and its current evidence
status.

\begin{table}[H]
  \centering
  \caption{System requirements and final evidence status.}
  \label{tab:requirements}
  \scriptsize
  \begin{tabularx}{\textwidth}{@{}p{0.08\textwidth}p{0.36\textwidth}p{0.18\textwidth}X@{}}
    \toprule
    ID & Requirement & Related question & Status boundary \\
    \midrule
    FR1 & Sentinel shall perform local initial visible-flame detection and
    expose the artifact identity and timing needed for analysis. & RQ1, RQ2 &
    Implemented and exercised on the stationary Sentinel. \\
    FR2 & A three-consecutive-frame temporal filter shall gate verification
    requests. & RQ3 & Implemented; its comparative effect requires the accepted
    ablation evidence. \\
    FR3 & Each verification request, Scout verdict and final decision shall be
    correlated by experiment, trial and event identity. & RQ3 & Implemented and
    exercised in stationary ROS~2 tests. \\
    FR4 & Scout shall perform independent full-frame visual verification and
    return at most one application verdict for an event. & RQ1, RQ3 &
    Implemented and exercised under bounded laboratory conditions. \\
    FR5 & Sentinel shall emit one terminal decision: \texttt{CONFIRMED},
    \texttt{REJECTED}, \texttt{TIMEOUT} or \texttt{ERROR}. & RQ3 & Implemented;
    physical alarm actuation is outside scope. \\
    FR6 & The mission layer shall model confirmed, rejected and degraded
    branches for both vehicles without enabling unverified physical flight. &
    RQ4 & Development simulation branches executed; the formal repeated
    campaign was not executed. \\
    NFR1 & Raw evidence shall be append-only and include configuration, timing,
    model identity and logging-drop counters. & All & Implemented in the enhanced
    protocol and experiment bundles. \\
    NFR2 & Runtime claims shall refer to an exact device, model hash, precision,
    input shape and service configuration. & RQ1, RQ2 & Mandatory analysis rule;
    historical 416-pixel and deployed 640-pixel Sentinel engines are distinct. \\
    NFR3 & A missing, duplicate, unrelated or late verdict shall not generate a
    second accepted decision. & RQ3, RQ4 & Enforced by the coordination state
    machine; bounded fault evidence is reported later. \\
    NFR4 & Evaluation shall keep offline, device, stationary, simulation and
    external evidence categories separate. & All & Thesis-wide reporting
    constraint. \\
    \bottomrule
  \end{tabularx}
\end{table}

No universal numeric latency, power or accuracy requirement was imposed after
the system had been measured. Instead, each experiment freezes an acceptance
criterion before its result is interpreted. This avoids converting an observed
value into a retrospective design target. Safety requirements are categorical:
physical integrated AI tests are stationary and propeller-free, and the
mission adapter must fail closed unless an exact simulator-only precondition is
satisfied.

\section{Sentinel and Scout platforms}
\label{sec:architecture:platforms}

The Sentinel is a customized Holybro X500 V2 platform with a Pixhawk~6C flight
controller, NVIDIA Jetson Nano 4~GB computer and Intel RealSense D435i camera.
Its role is broad-area hazard detection, temporal filtering, creation of target
metadata and publication of the final decision. NVIDIA's published processor
and memory specification provides resource context \cite{NvidiaJetsonNano};
the D400 family datasheet provides camera-family context
\cite{IntelD400Datasheet}. Actual camera mode, runtime and latency are taken
from experiment configuration and logs rather than inferred from those
specifications.

The Scout is a ModalAI PX4 Development Kit with a VOXL~2 computer based on the
QRB5165. Its role in the stationary workflow is to receive a correlated alert,
capture its own current image and perform full-frame secondary verification.
The official VOXL~2 feature matrix describes its processor, memory and
accelerators \cite{ModalAIVOXL2}. The retained runtime evidence, rather than
the feature matrix, determines which software backend was actually used.

\begin{table}[H]
  \centering
  \caption{Platform roles and architecture-relevant components. Specifications
  are nominal; measured configurations are reported with the experiments.}
  \label{tab:platforms}
  \small
  \begin{tabularx}{\textwidth}{@{}p{0.22\textwidth}XX@{}}
    \toprule
    Attribute & Sentinel & Scout \\
    \midrule
    Vehicle basis & Customized Holybro X500 V2 & ModalAI PX4 Development Kit \\
    Flight controller & Pixhawk 6C & PX4-compatible ModalAI platform \\
    Edge computer & Jetson Nano, 4~GB & VOXL~2 / QRB5165, 8~GB nominal memory \\
    Camera path & Intel RealSense D435i & Modal Pipe Architecture camera stream \\
    Perception role & Initial hazard detection and temporal filtering &
    Independent full-frame hazard verification \\
    Accepted deployment family & YOLOv8n / TensorRT & YOLOv8s / TensorFlow Lite \\
    Communication role & Alert publisher, verdict subscriber, decision
    publisher & Alert subscriber, verdict publisher \\
    Validated physical mode & Stationary, propellers removed & Stationary,
    propellers removed \\
    \bottomrule
  \end{tabularx}
\end{table}

\subsection{Physical configuration and scenario fit}

The two platforms are not interchangeable carriers for the same software. The
vehicles were assembled and configured for this project, with most physical
customization concentrated on the Sentinel. Its open X500 frame accommodated
the added Jetson Nano, D435i, wiring and power-conversion hardware, whereas the
Scout retained the more integrated ModalAI compute, camera and flight stack.
This difference supports the intended division between persistent initial
inspection and dispatched secondary verification, but it does not by itself
establish flight endurance, payload margin or maneuverability.

The Sentinel companion computer is supplied from the airframe power
distribution through an intermediate DC--DC regulator. A three-conductor
serial connection links the Jetson-side interface to the Pixhawk, and the
D435i is mechanically attached to and connected directly to the Jetson. The
accepted repository preserves the logical serial endpoint and baud rate, but
not a complete connector-level pin map, regulator part number, calibrated
voltage/current series or measured camera pose. The thesis therefore reports
the integration topology while treating electrical efficiency, energy use and
installed geometry as unmeasured quantities.

Table~\ref{tab:platform-constraint-evidence} separates resource constraints
that were measured from those that remain design context. This separation is
important: the term \emph{resource-constrained} is supported directly for
compute, memory, thermal observation, runtime compatibility and wireless
communication, but not by a calibrated electrical-power or airborne-endurance
experiment.

\begin{table}[H]
  \centering
  \caption{Platform constraints, evidence status and architectural consequence.}
  \label{tab:platform-constraint-evidence}
  \scriptsize
  \begin{tabularx}{\textwidth}{@{}p{0.17\textwidth}p{0.34\textwidth}X@{}}
    \toprule
    Constraint & Evidence in this study & Consequence for the architecture \\
    \midrule
    Compute and memory & Exact deployed runtimes, latency distributions, host
    CPU, available memory and swap were recorded on both edge computers. &
    Model, tensor layout and runtime are selected per platform rather than from
    predictive accuracy alone. \\
    Thermal operation & Host thermal-zone values were sampled during bounded
    stationary runs; cooling configuration and ambient conditions were not a
    controlled thermal experiment. & Sustained operation is evaluated directly,
    but no causal throttling or airborne cooling claim is made. \\
    Electrical power and energy & The Sentinel power topology is documented
    qualitatively; calibrated voltage, current, regulator efficiency and energy
    per decision were not measured. & No energy-efficiency ranking or flight-
    endurance improvement is claimed. \\
    Payload and geometry & Installed components are identified, but complete
    mass, rotor spacing, dimensions, centre of gravity and camera pose were not
    admitted as measured evidence. & Gazebo airframes remain provisional
    simulation proxies rather than independently identified physical models. \\
    Wireless communication & Compact ROS~2 request--reply events were measured
    over SwarmNet in stationary laboratory conditions. & Raw video remains
    local, reducing link load, while access-point coverage and interference
    remain deployment risks. \\
    Sensing & RGB camera paths were exercised on both vehicles; the final
    physical campaign did not validate depth, inertial, thermal, gas or
    particulate sensing for hazard decisions. & A second RGB decision improves
    traceability but is not independent physical confirmation of combustion. \\
    \bottomrule
  \end{tabularx}
\end{table}

The available sensor set is consequently broader than the evaluated sensing
path. The D435i family provides RGB, stereo depth and inertial capabilities
\cite{IntelD400Datasheet}, but the accepted Sentinel detector used only the
colour stream. Pixhawk telemetry was observed during the final endurance run,
yet its estimator and vehicle sensors were not inputs to the fire decision.
The Scout used its ModalAI camera stream through MPA; its PX4, IMU and QVIO
services were present as platform load but were not validated as parts of an
AI-controlled physical mission. Table~\ref{tab:sensor-use-boundary} makes this
availability--use distinction explicit.

\begin{table}[H]
  \centering
  \caption{Available platform capabilities versus their evaluated use.}
  \label{tab:sensor-use-boundary}
  \scriptsize
  \begin{tabularx}{\textwidth}{@{}p{0.23\textwidth}p{0.27\textwidth}X@{}}
    \toprule
    Component & Evaluated use & Boundary \\
    \midrule
    D435i colour camera & Sentinel RGB fire inference & Stationary live camera
    path exercised; no industrial-site calibration. \\
    D435i depth and IMU & None in the accepted hazard-decision path & Available
    device capabilities, not evaluated project results. \\
    Pixhawk 6C sensors and estimator & Receive-only telemetry/service
    coexistence & No integrated arming, navigation or physical mission. \\
    ModalAI camera through MPA & Scout current-frame RGB/YUV verification &
    Stationary bridge and decoding exercised, including unavailable/stale-frame
    faults. \\
    Scout PX4, IMU and QVIO services & Concurrent platform services during the
    endurance run & Their presence contributes load; navigation accuracy was
    not evaluated. \\
    \bottomrule
  \end{tabularx}
\end{table}

The labels ``recall-oriented'' and ``precision-oriented'' describe the intended
roles of the two stages. They are not assumed properties of YOLOv8n and
YOLOv8s. Operating thresholds and measured errors determine whether the final
configurations actually realize those priorities.

\section{Responsibility allocation and data flow}
\label{sec:architecture:dataflow}

Figure~\ref{fig:tiered-architecture} shows the stationary tiered architecture.
The Sentinel captures a frame, performs local inference and applies the
temporal filter. When the configured count is reached, it creates one alert
containing the event identity, detector confidence, temporal count, source
frame timing, model hash and normalized source-image box. ROS~2/DDS carries
this compact record over the SwarmNet wireless network.

The Scout treats the alert as a trigger. It does not crop its image around the
Sentinel box and does not navigate to the transferred coordinates. It performs
an independent full-frame inference on its current camera frame and publishes a
correlated \texttt{CONFIRMED}, \texttt{REJECTED} or \texttt{ERROR} verdict. The
Sentinel accepts only a verdict for the active event and publishes one terminal
decision. A timeout is generated locally if a matching verdict does not arrive
within the configured interval.

\begin{figure}[H]
  \centering
  \resizebox{\textwidth}{!}{\begin{tikzpicture}[
  font=\small,
  block/.style={draw,rounded corners=2pt,align=center,text width=34mm,
    minimum height=13mm,inner sep=2.2mm,fill=blue!4},
  decision/.style={block,fill=orange!10},
  flow/.style={-{Latex[length=2.2mm]},thick},
  exchange/.style={-{Latex[length=2.2mm]},thick,blue!65!black},
  future/.style={-{Latex[length=2.2mm]},thick,dashed,gray!75},
  label/.style={font=\scriptsize,align=center,fill=white,inner sep=1.5pt}
]
  % The enlarged device boundaries deliberately provide a visible margin around
  % every processing block after the complete figure is scaled to text width.
  \draw[rounded corners=4pt,fill=blue!1] (-2.10,0.30) rectangle (16.15,3.35);
  \draw[rounded corners=4pt,fill=blue!1] (-2.10,-3.45) rectangle (16.15,-0.30);
  \node[font=\bfseries\small,fill=blue!1,inner sep=2pt] at (7.05,3.00)
    {Sentinel --- Jetson Nano};
  \node[font=\bfseries\small,fill=blue!1,inner sep=2pt] at (7.05,-3.08)
    {Scout --- VOXL 2};

  \node[block] (scam) at (0,1.65) {D435i camera\\local image stream};
  \node[block] (sdet) at (4.70,1.65) {YOLOv8n / TensorRT\\initial detection};
  \node[block] (filter) at (9.40,1.65) {three-frame\\temporal filter};
  \node[decision] (coord) at (14.10,1.65) {event coordinator\\final decision};

  \node[block] (ccam) at (0,-1.70) {Scout camera\\local image stream};
  \node[block] (buffer) at (4.70,-1.70) {latest-frame buffer\\freshness gate};
  \node[block] (cdet) at (9.40,-1.70) {YOLOv8s / TFLite\\full-frame verification};
  \node[decision] (verdict) at (14.10,-1.70) {correlated verdict\\for active event};

  \draw[flow] (scam) -- (sdet);
  \draw[flow] (sdet) -- (filter);
  \draw[flow] (filter) -- (coord);
  \draw[flow] (ccam) -- (buffer);
  \draw[flow] (buffer) -- (cdet);
  \draw[flow] (cdet) -- (verdict);

  \draw[exchange] ([xshift=-4mm]coord.south) --
    node[label,left=2mm]{\texttt{HazardAlert}\\event identity + metadata}
    ([xshift=-4mm]verdict.north);
  \draw[exchange] ([xshift=4mm]verdict.north) --
    node[label,right=2mm]{\texttt{HazardVerdict}\\confirm, reject or error}
    ([xshift=4mm]coord.south);

  \node[decision,below=18mm of verdict] (log) {terminal decision\\and persistent log};
  \node[block,fill=gray!8,left=11mm of log] (flight) {flight-control adapter\\dry-run / simulation only};
  \draw[flow] (verdict.south) -- (log.north);
  \draw[future] (log.west) -- (flight.east);
\end{tikzpicture}
}
  \caption{Tiered edge architecture. Solid paths belong to the evaluated
  stationary perception and ROS~2 workflow. The dashed flight-control path is
  implemented for constrained dry-run and simulation use but was not validated
  as integrated physical flight.}
  \label{fig:tiered-architecture}
\end{figure}

This division keeps raw video local to each device and limits the network
exchange to event metadata and decisions. It also means that verification is
possible only when the Scout's current field of view contains relevant visual
evidence. The tiered workflow provides no implicit camera-to-camera correspondence.
That constraint is part of the architecture rather than a later experimental
failure.

\subsection{Onboard execution and ground-station boundary}

The physical perception architecture is fully onboard in the computational
sense. Sentinel camera capture, TensorRT inference, temporal filtering, alert
publication and final-decision handling execute on the Jetson Nano; Scout
camera acquisition, decoding and TFLite verification execute on VOXL~2. The
Sentinel also hosts the SwarmNet wireless access point. The laboratory laptop
was used for configuration, bounded test orchestration, monitoring and
evidence retrieval, but it was not the inference server or the required final-
decision processor in the evaluated two-device path. Loss of an external
Internet or cloud service therefore does not, by design, remove local
perception.

This property must not be described as a \emph{full aerial experiment}. The
accepted physical evidence is stationary and propeller-free; aerial state
transitions were exercised only in simulation. A precise description is a
\emph{fully onboard edge-processing architecture with stationary physical
validation and simulated aerial mission execution}.

Keeping images local reduces radio traffic and avoids making hazard decisions
dependent on continuous raw-video transport. It also distributes failure and
maintenance across two dissimilar embedded environments. The Sentinel access
point and coordinator concentrate connectivity and finalization responsibility,
so their loss can isolate the Scout even though inference is decentralized.
Within a warehouse-like scenario, local processing is attractive where
external connectivity is weak or restricted, but metallic structures, aisles,
occlusion, multipath propagation and changing illumination can degrade both
radio and RGB performance. These are design implications, not measured
industrial-site results.

\section{Communication interfaces and coordinate conventions}
\label{sec:architecture:interfaces}

The enhanced interface uses three typed ROS~2 messages. Their roles are
summarized in Table~\ref{tab:interfaces}. All messages carry a ROS header and
the identifiers needed to join records across the workflow.

\begin{table}[H]
  \centering
  \caption{Enhanced decision interfaces.}
  \label{tab:interfaces}
  \scriptsize
  \begin{tabularx}{\textwidth}{@{}p{0.23\textwidth}p{0.17\textwidth}X@{}}
    \toprule
    Topic and type & Direction & Principal semantics \\
    \midrule
    \path{/swarm/hazard_alert} / \texttt{HazardAlert} & Sentinel to Scout &
    Experiment, trial and event IDs; source device; hazard type; confidence;
    temporal count; source/alert monotonic times; model hash; normalized
    source-image box and frame dimensions. \\
    \path{/swarm/hazard_verdict} / \texttt{HazardVerdict} & Scout to Sentinel &
    Matching IDs; verdict; hazard type and confidence; receipt/source/verdict
    times; Scout model hash; reason. \\
    \path{/swarm/hazard_decision} / \texttt{HazardDecision} & Sentinel to
    logger or consumer & Matching IDs; terminal decision; both confidences;
    detection, alert, verdict and decision times; reason. \\
    \bottomrule
  \end{tabularx}
\end{table}

The alert fields \texttt{target\_u}, \texttt{target\_v},
\texttt{bbox\_width} and \texttt{bbox\_height} are normalized in the Sentinel
source-camera frame. The origin is the top-left; \(u\) increases to the right
and \(v\) increases downwards. These values are not latitude, longitude, local
vehicle coordinates or Scout image coordinates. Earlier legacy
\texttt{TargetHazard} messages used ambiguous location naming. The enhanced
message makes the image-space meaning explicit and is the normative interface
for the final evaluation.

ROS header time supports cross-device record alignment when clock
synchronization is known. Each \texttt{*\_monotonic\_ns} field is valid for
durations on its originating device. Monotonic values from different devices
cannot be subtracted directly. Application round-trip time can be calculated
on the initiating Sentinel without assuming that the two monotonic clock
epochs are equal.

Figure~\ref{fig:decision-sequence} makes the request, response and timeout paths
explicit. Event identifiers are opaque and cannot be reused within an
experiment. A duplicate, late or unrelated verdict cannot close a different
active event.

\begin{figure}[p]
  \centering
  \resizebox{0.96\textwidth}{!}{\begin{tikzpicture}[
  font=\small,
  actor/.style={draw,rounded corners=2pt,minimum width=30mm,minimum height=9mm,
    fill=blue!5,font=\bfseries\small,align=center},
  flow/.style={-{Latex[length=2mm]},thick},
  timeout/.style={-{Latex[length=2mm]},thick,dashed,gray!75},
  msg/.style={font=\scriptsize,align=center,fill=white,inner sep=1.2pt,
    text width=37mm},
  action/.style={draw,rounded corners=2pt,fill=orange!8,font=\scriptsize,
    align=center,text width=31mm,minimum height=8mm,inner sep=1.5mm}
]
  \node[actor] (perception) at (0,0) {Sentinel\\perception};
  \node[actor] (coordinator) at (4.75,0) {Sentinel\\coordinator};
  \node[actor] (dds) at (9.50,0) {ROS 2 / DDS};
  \node[actor] (scout) at (14.25,0) {Scout\\verifier};
  \foreach \x in {0,4.75,9.50,14.25}
    \draw[densely dashed,gray!70] (\x,-0.48) -- (\x,-16.45);

  \node[action] at (0,-1.40) {D435i frame $\rightarrow$ TensorRT\\detection and decoding};
  \draw[flow] (0,-2.80) --
    node[msg,above]{candidate: class, confidence, box and source time}
    (4.75,-2.80);

  \node[action] at (4.75,-4.15) {apply three-frame rule; create\\unique event; start deadline};
  \draw[flow] (4.75,-5.55) --
    node[msg,above]{\texttt{HazardAlert}: experiment, trial, event and model IDs}
    (9.50,-5.55);
  \draw[flow] (9.50,-6.80) --
    node[msg,above]{deliver alert with normalized image metadata}
    (14.25,-6.80);

  \node[action] at (14.25,-8.20) {validate IDs, freshness and duplicate state;\\snapshot own frame; full-frame inference};
  \draw[flow] (14.25,-10.20) --
    node[msg,above]{\texttt{HazardVerdict}: matching event, decision, score and reason}
    (9.50,-10.20);
  \draw[flow] (9.50,-11.55) --
    node[msg,above]{deliver correlated confirmation, rejection or error}
    (4.75,-11.55);

  \node[action] at (4.75,-13.05) {require the active event ID; ignore stale or\\duplicate verdicts; persist one terminal record};
  \draw[flow] (4.75,-14.65) --
    node[msg,above]{publish exactly one \texttt{HazardDecision}; reset}
    (9.50,-14.65);
  \draw[timeout] (9.50,-16.05) --
    node[msg,above]{alternative: no matching verdict before deadline}
    (4.75,-16.05);
\end{tikzpicture}
}
  \caption{Detailed correlated Sentinel--Scout decision sequence. The Scout
  validates the request and performs independent full-frame inference on its
  own fresh image. The coordinator accepts only a matching active-event
  verdict and records exactly one terminal outcome; the dashed path is the
  deadline-expiry alternative. Image-space target metadata is retained for
  traceability and is not used to steer the Scout in the evaluated stationary
  workflow.}
  \label{fig:decision-sequence}
\end{figure}
\clearpage

\section{Coordination states and mission boundary}
\label{sec:architecture:states}

The perception coordinator and the flight mission controller use related but
distinct state models. In the stationary workflow, the coordinator begins in
\texttt{SEARCHING}. A qualifying candidate moves it to \texttt{TRACKING}; an
alert moves it to \texttt{AWAITING\_SCOUT}. A matching verdict produces
\texttt{CONFIRMED}, \texttt{REJECTED} or \texttt{ERROR}; expiration produces
\texttt{TIMEOUT}. Exactly one terminal decision is recorded before the
coordinator resets to \texttt{SEARCHING}.

\begin{figure}[H]
  \centering
  \resizebox{\textwidth}{!}{\begin{tikzpicture}[
  state/.style={draw,rounded corners=2pt,minimum width=27mm,minimum height=10mm,
    align=center,font=\small,fill=blue!4},
  terminal/.style={state,fill=orange!10},
  reset/.style={state,fill=gray!10},
  flow/.style={-{Latex[length=2mm]},thick},
  return/.style={-{Latex[length=2mm]},thick,gray!70},
  edge/.style={font=\scriptsize,align=center,fill=white,inner sep=1.5pt}
]
  \node[state] (search) at (0,0) {SEARCHING};
  \node[state] (track) at (5.30,0) {TRACKING};
  \node[state] (await) at (10.60,0) {AWAITING\\SCOUT};

  \draw[flow] (search.east) -- (track.west);
  \node[edge] at (2.65,0.72) {qualifying candidate};
  \draw[flow] (track.east) -- (await.west);
  \node[edge] at (7.95,0.78) {three-frame trigger;\\publish alert};

  \node[terminal] (confirm) at (3.10,-2.40) {CONFIRMED\\matching verdict};
  \node[terminal] (reject) at (8.25,-2.40) {REJECTED\\matching verdict};
  \node[terminal] (fail) at (13.40,-2.40) {TIMEOUT / ERROR\\deadline or fault};
  \node[reset] (record) at (8.25,-4.70) {record one terminal\\decision and reset};

  \draw[flow] (await.south west) -- (confirm.north);
  \draw[flow] (await.south) -- (reject.north);
  \draw[flow] (await.south east) -- (fail.north);

  \draw[return] (confirm.south) |- (record.west);
  \draw[return] (reject.south) -- (record.north);
  \draw[return] (fail.south) |- (record.east);
  \draw[return] (record.west) -- ++(-8.25,0) -|
    node[edge,near start,below]{next event} (search.south);
\end{tikzpicture}
}
  \caption{Stationary perception and coordination states. Terminal decisions
  are recorded once before the coordinator returns to search.}
  \label{fig:coordination-states}
\end{figure}

The mission layer adds vehicle states such as \path{PREFLIGHT},
\path{TAKEOFF}, \path{PATROL}, \path{SUSPECT_HOLD},
\path{SCOUT_DISPATCH}, \path{VERIFYING},
\path{CONFIRMED_HOLD}, \path{REJECTED_RESUME}, \path{DEGRADED},
\path{RTL} and \path{LANDED}. Its target positions are scenario-defined
local-ENU coordinates. They are not derived from the Sentinel image box. Two
development PX4/Gazebo runs exercised the nominal confirmed and rejected
branches on provisional vehicle proxies. A later accepted visual run executed
the confirmed branch from preflight through patrol, hold, Scout dispatch,
verification, alarm publication, return and landing while retaining the
corresponding state-transition log. These runs establish executable control
paths under their recorded simulation conditions, not a reliability rate or
physical-flight validation.

The formal repeated simulation campaign was not executed before the final
evidence freeze. Measured physical parameters and the complete preregistered
fault set did not close every readiness gate. Consequently, RQ4 receives only
a bounded answer from the retained development runs. The visual record is
supporting evidence for one run and does not replace the formal repeated
campaign or its machine-readable telemetry.

\section{Resource, logging and safety constraints}
\label{sec:architecture:constraints}

The Sentinel has less nominal memory and an older GPU architecture than the
Scout platform, but hardware specifications alone do not determine the model
assignment. Runtime availability and graph compatibility also matter. The
Sentinel therefore uses a TensorRT deployment path, while the accepted Scout
configuration uses TensorFlow Lite. Candidate screening remains broader than
these historical deployments so that the thesis can distinguish a deployed
artifact from a model selected under the later controlled comparison.

Logging is moved off the inference path through a bounded queue and a dedicated
writer. This reduces synchronous file-I/O interference while making overload
visible through drop counters. Raw JSON Lines records are append-only.
Corrections, exclusions and derived metrics are stored separately so that an
analysis cannot silently replace the observation from which it was produced.

The communication layer uses explicit correlation and a timeout rather than
equating DDS reliability with successful application completion. A time-limited
rejection cooldown can prevent immediate repeated requests for the same
normalized image region. Persistent evidence or a sufficient confidence
increase can override the cooldown, and expiry makes the region eligible
again. This policy must be evaluated; it is not assumed to improve detection.

The physical safety boundary is strict. Integrated perception experiments use
stationary aircraft with removed propellers. The architecture does not convert
an image coordinate into a flight setpoint. The flight adapter is used only
under explicit dry-run or simulator safeguards. Physical arming, autonomous
movement and industrial operation require separate validation outside the
evidence presented here.

\section{Design rationale and research-question traceability}
\label{sec:architecture:rationale}

The architecture addresses RQ1 by allowing model and runtime choices to be
evaluated per device rather than assuming homogeneous hardware. It addresses
RQ2 by defining the concurrent camera, inference, communication and logging
services whose sustained operation must be measured. It addresses RQ3 through
explicit single-frame, temporal-filter and verification stages with measurable
timestamps and terminal outcomes. The separate mission controller supplies the
software boundary for RQ4, while keeping simulator coordinates and physical
perception evidence distinct. RQ5 remains principally a data and evaluation
question; the architecture contributes the frozen model and preprocessing
contracts needed to test external transfer.

The expected benefit is fewer unsupported terminal confirmations than a
single-frame Sentinel decision, at the cost of additional delay, energy use and
failure opportunities. That statement is a hypothesis. The architecture makes
the comparison observable, but only the accepted experiments in Chapters~6
and 7 can determine its magnitude under defined conditions.
